\subsubsection{$q$-binomial formulas}

The properties we have seen so far are suggesting that $q$-binomial
coefficients not only generalize binomial coefficients, but also share most of
their properties in a somewhat modified form. In other words, we start
expecting most properties of binomial coefficients to generalize to
$q$-binomial coefficients, often in several ways (e.g., the recurrence of the
binomial coefficients generalized in two ways).

Let us see how this expectation holds up for the most famous property of
binomial coefficients: the binomial formula%
\[
\left(  a+b\right)  ^{n}=\sum_{k=0}^{n}\dbinom{n}{k}a^{k}b^{n-k}.
\]
This formula holds whenever $a$ and $b$ are two elements of a commutative
ring, or even more generally, whenever $a$ and $b$ are two commuting elements
of an arbitrary ring. If we want to integrate a $q$ into this formula, we need to

\begin{itemize}
\item either change the structure of the formula,

\item or modify the commutativity assumption.
\end{itemize}

This gives rise to two different \textquotedblleft$q$%
-analogues\textquotedblright\ of the binomial formula. Both are important (one
for the theory of partitions, and another for the theory of
\href{https://www.ams.org/notices/200601/what-is.pdf}{quantum groups}). Here
is the first one:

\begin{theorem}
[1st $q$-binomial theorem]\label{thm.pars.qbinom.binom1}Let $K$ be a
commutative ring. Let $a,b\in K$ and $n\in\mathbb{N}$. In the polynomial ring
$K\left[  q\right]  $, we have%
\[
\left(  aq^{0}+b\right)  \left(  aq^{1}+b\right)  \cdots\left(  aq^{n-1}%
+b\right)  =\sum_{k=0}^{n}q^{k\left(  k-1\right)  /2}\dbinom{n}{k}_{q}%
a^{k}b^{n-k}.
\]

\end{theorem}

Note that setting $q=1$ in Theorem \ref{thm.pars.qbinom.binom1} (i.e.,
substituting $1$ for $q$) recovers the good old binomial formula, since all
the $n$ factors on the left hand side become $a+b$.

There is a straightforward way to prove Theorem \ref{thm.pars.qbinom.binom1}
by induction on $n$ (see Exercise \ref{exe.pars.qbinom.qbinom} \textbf{(a)}).
Let us instead give a nicer argument. This argument will rely on the following
general fact:

\begin{lemma}
\label{lem.prodrule.sum-ai-plus-bi}Let $L$ be a commutative ring. Let
$n\in\mathbb{N}$. Let $\left[  n\right]  $ denote the set $\left\{
1,2,\ldots,n\right\}  $. Let $a_{1},a_{2},\ldots,a_{n}$ be $n$ elements of
$L$. Let $b_{1},b_{2},\ldots,b_{n}$ be $n$ further elements of $L$. Then,%
\begin{equation}
\prod_{i=1}^{n}\left(  a_{i}+b_{i}\right)  =\sum_{S\subseteq\left[  n\right]
}\left(  \prod_{i\in S}a_{i}\right)  \left(  \prod_{i\in\left[  n\right]
\setminus S}b_{i}\right)  . \label{eq.lem.prodrule.sum-ai-plus-bi.eq}%
\end{equation}

\end{lemma}

Lemma \ref{lem.prodrule.sum-ai-plus-bi} is well-known and intuitively clear:
When expanding the product $\prod_{i=1}^{n}\left(  a_{i}+b_{i}\right)
=\left(  a_{1}+b_{1}\right)  \left(  a_{2}+b_{2}\right)  \cdots\left(
a_{n}+b_{n}\right)  $, you obtain a sum of $2^{n}$ terms, each of which is a
product of one addend chosen from each of the $n$ sums $a_{1}+b_{1}%
,a_{2}+b_{2},\ldots,a_{n}+b_{n}$. This is precisely what the right hand side
of (\ref{eq.lem.prodrule.sum-ai-plus-bi.eq}) is. A rigorous proof of Lemma
\ref{lem.prodrule.sum-ai-plus-bi} can be found in \cite[Exercise 6.1
\textbf{(a)}]{detnotes}.

\begin{proof}
[Proof of Theorem \ref{thm.pars.qbinom.binom1}.]Let $\left[  n\right]  $
denote the set $\left\{  1,2,\ldots,n\right\}  $. We have%
\begin{align*}
&  \left(  aq^{0}+b\right)  \left(  aq^{1}+b\right)  \cdots\left(
aq^{n-1}+b\right)  =\prod_{i=1}^{n}\left(  aq^{i-1}+b\right) \\
&  =\sum_{S\subseteq\left[  n\right]  }\underbrace{\left(  \prod_{i\in
S}\left(  aq^{i-1}\right)  \right)  }_{=a^{\left\vert S\right\vert }%
\prod_{i\in S}q^{i-1}}\underbrace{\left(  \prod_{i\in\left[  n\right]
\setminus S}b\right)  }_{=b^{\left\vert \left[  n\right]  \setminus
S\right\vert }}\\
&  \ \ \ \ \ \ \ \ \ \ \ \ \ \ \ \ \ \ \ \ \left(  \text{by Lemma
\ref{lem.prodrule.sum-ai-plus-bi}, applied to }L=K\left[  q\right]  \text{,
}a_{i}=aq^{i-1}\text{ and }b_{i}=b\right) \\
&  =\underbrace{\sum_{S\subseteq\left[  n\right]  }}_{=\sum\limits_{k=0}%
^{n}\ \ \sum\limits_{\substack{S\subseteq\left[  n\right]  ;\\\left\vert
S\right\vert =k}}}a^{\left\vert S\right\vert }\underbrace{\left(  \prod_{i\in
S}q^{i-1}\right)  }_{\substack{=q^{\operatorname*{sum}S-\left\vert
S\right\vert }\\\text{(since the sum of the exponents }i-1\\\text{over all
}i\in S\text{ is precisely }\operatorname*{sum}S-\left\vert S\right\vert
\text{)}}}b^{\left\vert \left[  n\right]  \setminus S\right\vert }\\
&  =\sum_{k=0}^{n}\ \ \sum_{\substack{S\subseteq\left[  n\right]
;\\\left\vert S\right\vert =k}}\ \ \underbrace{a^{\left\vert S\right\vert }%
}_{\substack{=a^{k}\\\text{(since }\left\vert S\right\vert =k\text{)}%
}}\ \ \underbrace{q^{\operatorname*{sum}S-\left\vert S\right\vert }%
}_{\substack{=q^{\operatorname*{sum}S-k}\\\text{(since }\left\vert
S\right\vert =k\text{)}}}\ \ \underbrace{b^{\left\vert \left[  n\right]
\setminus S\right\vert }}_{\substack{=b^{n-k}\\\text{(since }S\text{ is a
}k\text{-element}\\\text{subset of the }n\text{-element}\\\text{set }\left[
n\right]  \text{, and thus we}\\\text{have }\left\vert \left[  n\right]
\setminus S\right\vert =n-k\text{)}}}\\
&  =\sum_{k=0}^{n}\ \ \sum_{\substack{S\subseteq\left[  n\right]
;\\\left\vert S\right\vert =k}}a^{k}\underbrace{q^{\operatorname*{sum}S-k}%
}_{\substack{=q^{\operatorname*{sum}S-\left(  1+2+\cdots+k\right)
}q^{1+2+\cdots+\left(  k-1\right)  }\\=q^{\operatorname*{sum}S-\left(
1+2+\cdots+k\right)  }q^{k\left(  k-1\right)  /2}\\\text{(since }%
1+2+\cdots+\left(  k-1\right)  =k\left(  k-1\right)  /2\text{)}}}b^{n-k}\\
&  =\sum_{k=0}^{n}\ \ \sum_{\substack{S\subseteq\left[  n\right]
;\\\left\vert S\right\vert =k}}a^{k}q^{\operatorname*{sum}S-\left(
1+2+\cdots+k\right)  }q^{k\left(  k-1\right)  /2}b^{n-k}\\
&  =\sum_{k=0}^{n}q^{k\left(  k-1\right)  /2}\underbrace{\left(
\sum\limits_{\substack{S\subseteq\left[  n\right]  ;\\\left\vert S\right\vert
=k}}q^{\operatorname*{sum}S-\left(  1+2+\cdots+k\right)  }\right)
}_{\substack{=\dbinom{n}{k}_{q}\\\text{(by Proposition
\ref{prop.pars.qbinom.alt-defs} \textbf{(b)},}\\\text{since }\left[  n\right]
=\left\{  1,2,\ldots,n\right\}  \text{)}}}a^{k}b^{n-k}=\sum_{k=0}%
^{n}q^{k\left(  k-1\right)  /2}\dbinom{n}{k}_{q}a^{k}b^{n-k}.
\end{align*}
This proves Theorem \ref{thm.pars.qbinom.binom1}.
\end{proof}

The 2nd $q$-binomial theorem grows out of noncommutativity:

\begin{theorem}
[2nd $q$-binomial theorem, aka Potter's binomial theorem]%
\label{thm.pars.qbinom.binom2}Let $L$ be a commutative ring. Let $\omega\in
L$. Let $A$ be a noncommutative $L$-algebra. Let $a,b\in A$ be such that
$ba=\omega ab$. Then,%
\[
\left(  a+b\right)  ^{n}=\sum_{k=0}^{n}\dbinom{n}{k}_{\omega}a^{k}b^{n-k}.
\]

\end{theorem}

The condition $ba=\omega ab$ looks somewhat artificial -- do such elements
$a,b$ actually exist in the wild? Indeed they do, as the following examples show:

\begin{example}
Let $L=\mathbb{Z}$ and $\omega=-1$ and $A=\mathbb{Z}^{2\times2}$ (the ring of
$2\times2$-matrices with integer entries). Let%
\[
a=\left(
\begin{array}
[c]{cc}%
0 & 1\\
1 & 0
\end{array}
\right)  \ \ \ \ \ \ \ \ \ \ \text{and}\ \ \ \ \ \ \ \ \ \ b=\left(
\begin{array}
[c]{cc}%
1 & 0\\
0 & -1
\end{array}
\right)  .
\]
It is easy to check that these two matrices satisfy $ba=-ab$, that is,
$ba=\omega ab$. Thus, Theorem \ref{thm.pars.qbinom.binom2} predicts that%
\[
\left(  a+b\right)  ^{n}=\sum_{k=0}^{n}\dbinom{n}{k}_{\omega}a^{k}b^{n-k}.
\]
And this is indeed true (check it for $n=3$).
\end{example}

\begin{example}
Let $L=\mathbb{R}$. Let $A$ be the ring of $\mathbb{R}$-linear operators on
$C^{\infty}\left(  \mathbb{R}\right)  =\left\{  \text{smooth functions from
}\mathbb{R}\text{ to }\mathbb{R}\right\}  $. Let $\omega$ be any real number.

Let $b\in A$ be the differentiation operator (sending each $f\in C^{\infty
}\left(  \mathbb{R}\right)  $ to $f^{\prime}$).

Let $a\in A$ be the operator that substitutes $\omega x$ for $x$ in the
function (in other words, it shrinks the plot of the function by $\omega$ in
the $x$-direction).

Then, you can check that $ba=\omega ab$. (Indeed, $\left(  f\left(  \omega
x\right)  \right)  ^{\prime}=\omega f^{\prime}\left(  \omega x\right)  $.)
\end{example}

The proof of Theorem \ref{thm.pars.qbinom.binom2} is again a homework exercise
(Exercise \ref{exe.pars.qbinom.qbinom} \textbf{(b)}).

The two $q$-binomial theorems are not entirely unrelated: Theorem
\ref{thm.pars.qbinom.binom1} can be obtained from Theorem
\ref{thm.pars.qbinom.binom2}. (See Exercise \ref{exe.pars.qbinom.2to1} for the details.)

\subsubsection{Counting subspaces of vector spaces}

We have introduced the $q$-binomial coefficient $\dbinom{n}{k}_{q}$ as a
generating function for a certain sort of partitions -- i.e., a
\textquotedblleft weighted number\textquotedblright\ of partitions, where each
partition $\lambda$ has weight $q^{\left\vert \lambda\right\vert }$. However,
for certain integers $a$, the number $\dbinom{n}{k}_{a}$ has other
interpretations, too. A particularly striking one can be found when $a$ is the
size of a finite field.

Let us recall a few things about finite fields (see, e.g., \cite[Chapter
5]{23wa} or \cite[Theorem 15.7.3]{Artin}):

\begin{itemize}
\item For any prime power $p^{k}$, there is a finite field of size $p^{k}$; it
is unique up to isomorphism, and is therefore often called the
\textquotedblleft Galois field of size $p^{k}$\textquotedblright\ and denoted
by $\mathbb{F}_{p^{k}}$. The finite fields $\mathbb{F}_{p^{1}}$ are easiest to
construct -- they are just the quotient rings $\mathbb{Z}/p\mathbb{Z}%
=\mathbb{Z}/p$ (that is, the rings of integers modulo $p$). Higher prime
powers are more complicated. For example, the finite field $\mathbb{F}_{p^{2}%
}$ can be obtained by starting with $\mathbb{Z}/p$ and adjoining a square root
of an element that is not a square. It is not $\mathbb{Z}/p^{2}$, since
$\mathbb{Z}/p^{2}$ is not a field!

\item Linear algebra (i.e., the notions of vector spaces, subspaces, linear
independence, bases, matrices, Gaussian elimination, etc.) can be done over
any field. In fact, many of its concepts can be defined over any commutative
ring, but only over fields do they behave as nicely as they do over the real
numbers. Thus, much of the linear algebra that you have learned over the real
numbers remains valid over any field. (Exceptions are some properties that
rely on positivity or on characteristic $0$.)
\end{itemize}

Thus, it makes sense to talk about finite-dimensional vector spaces over
finite fields. Such spaces are finite as sets, and thus can be viewed as
combinatorial objects. An $n$-dimensional vector space over a finite field $F$
has size $\left\vert F\right\vert ^{n}$.

Now, we might wonder how many $k$-dimensional subspaces such an $n$%
-dimensional vector space has. The answer is given by the following theorem:

\begin{theorem}
\label{thm.pars.qbinom.subsp-count}Let $F$ be a finite field. Let
$n,k\in\mathbb{N}$. Let $V$ be an $n$-dimensional $F$-vector space. Then,%
\[
\dbinom{n}{k}_{\left\vert F\right\vert }=\left(  \text{\# of }%
k\text{-dimensional vector subspaces of }V\right)  .
\]

\end{theorem}

Compare this with the classical fact that if $S$ is an $n$-element set, then%
\[
\dbinom{n}{k}=\left(  \text{\# of }k\text{-element subsets of }S\right)  .
\]
This hints at an analogy between finite sets and finite-dimensional vector
spaces. Such an analogy does indeed exist; the expository paper \cite{Cohn04}
gives a great overview over its reach.

The easiest proof of Theorem \ref{thm.pars.qbinom.subsp-count} uses three
lemmas. The first one is a classical fact from linear algebra, which holds for
any vector space (not necessarily finite-dimensional) over any field (not
necessarily finite):

\begin{lemma}
\label{lem.linalg.lin-ind-via-span}Let $F$ be a field. Let $V$ be an
$F$-vector space. Let $\left(  v_{1},v_{2},\ldots,v_{k}\right)  $ be a
$k$-tuple of vectors in $V$. Then, $\left(  v_{1},v_{2},\ldots,v_{k}\right)  $
is linearly independent if and only if each $i\in\left\{  1,2,\ldots
,k\right\}  $ satisfies $v_{i}\notin\operatorname*{span}\left(  v_{1}%
,v_{2},\ldots,v_{i-1}\right)  $ (where the span $\operatorname*{span}\left(
{}\right)  $ of an empty list is understood to be the set $\left\{
\mathbf{0}\right\}  $ consisting only of the zero vector $\mathbf{0}$). In
other words, $\left(  v_{1},v_{2},\ldots,v_{k}\right)  $ is linearly
independent if and only if we have%
\begin{align*}
v_{1}  &  \notin\underbrace{\operatorname*{span}\left(  {}\right)
}_{=\left\{  \mathbf{0}\right\}  }\ \ \ \ \ \ \ \ \ \ \text{and}\\
v_{2}  &  \notin\operatorname*{span}\left(  v_{1}\right)
\ \ \ \ \ \ \ \ \ \ \text{and}\\
v_{3}  &  \notin\operatorname*{span}\left(  v_{1},v_{2}\right)
\ \ \ \ \ \ \ \ \ \ \text{and}\\
&  \cdots\ \ \ \ \ \ \ \ \ \ \text{and}\\
v_{k}  &  \notin\operatorname*{span}\left(  v_{1},v_{2},\ldots,v_{k-1}\right)
.
\end{align*}

\end{lemma}

\begin{proof}
[Proof of Lemma \ref{lem.linalg.lin-ind-via-span}.]We must prove that $\left(
v_{1},v_{2},\ldots,v_{k}\right)  $ is linearly independent if and only if each
$i\in\left\{  1,2,\ldots,k\right\}  $ satisfies $v_{i}\notin%
\operatorname*{span}\left(  v_{1},v_{2},\ldots,v_{i-1}\right)  $. This is an
\textquotedblleft if and only if\textquotedblright\ statement; we shall prove
its \textquotedblleft only if\textquotedblright\ (i.e., \textquotedblleft%
$\Longrightarrow$\textquotedblright) and \textquotedblleft
if\textquotedblright\ (i.e., \textquotedblleft$\Longleftarrow$%
\textquotedblright) directions separately:

\begin{enumerate}
\item[$\Longrightarrow:$] Assume that the $k$-tuple $\left(  v_{1}%
,v_{2},\ldots,v_{k}\right)  $ is linearly independent. Let $i\in\left\{
1,2,\ldots,k\right\}  $. If we had $v_{i}\in\operatorname*{span}\left(
v_{1},v_{2},\ldots,v_{i-1}\right)  $, then we could write $v_{i}$ in the form
$v_{i}=\alpha_{1}v_{1}+\alpha_{2}v_{2}+\cdots+\alpha_{i-1}v_{i-1}$ for some
coefficients $\alpha_{1},\alpha_{2},\ldots,\alpha_{i-1}\in F$, and therefore
these coefficients $\alpha_{1},\alpha_{2},\ldots,\alpha_{i-1}$ would satisfy
\begin{align*}
&  \underbrace{\alpha_{1}v_{1}+\alpha_{2}v_{2}+\cdots+\alpha_{i-1}v_{i-1}%
}_{=v_{i}}+\underbrace{\left(  -1\right)  v_{i}}_{=-v_{i}}%
+\underbrace{0v_{i+1}+0v_{i+2}+\cdots+0v_{k}}_{=\mathbf{0}}\\
&  =v_{i}+\left(  -v_{i}\right)  +\mathbf{0}=\mathbf{0},
\end{align*}
which would be a nontrivial linear dependence relation between $\left(
v_{1},v_{2},\ldots,v_{k}\right)  $ (nontrivial because $v_{i}$ appears in it
with the nonzero coefficient $-1$); this would contradict the linear
independence of $\left(  v_{1},v_{2},\ldots,v_{k}\right)  $. Hence, we cannot
have $v_{i}\in\operatorname*{span}\left(  v_{1},v_{2},\ldots,v_{i-1}\right)
$. In other words, we have $v_{i}\notin\operatorname*{span}\left(  v_{1}%
,v_{2},\ldots,v_{i-1}\right)  $.

Forget that we fixed $i$. We thus have shown that each $i\in\left\{
1,2,\ldots,k\right\}  $ satisfies $v_{i}\notin\operatorname*{span}\left(
v_{1},v_{2},\ldots,v_{i-1}\right)  $. This proves the \textquotedblleft%
$\Longrightarrow$\textquotedblright\ direction of our claim.

\item[$\Longleftarrow:$] Assume that each $i\in\left\{  1,2,\ldots,k\right\}
$ satisfies $v_{i}\notin\operatorname*{span}\left(  v_{1},v_{2},\ldots
,v_{i-1}\right)  $. We must prove that the $k$-tuple $\left(  v_{1}%
,v_{2},\ldots,v_{k}\right)  $ is linearly independent. Indeed, assume the
contrary. Thus, this $k$-tuple is linearly dependent. In other words, there
exist coefficients $\beta_{1},\beta_{2},\ldots,\beta_{k}\in F$ that satisfy
$\beta_{1}v_{1}+\beta_{2}v_{2}+\cdots+\beta_{k}v_{k}=\mathbf{0}$ and that are
not all zero. Consider these $\beta_{1},\beta_{2},\ldots,\beta_{k}$. At least
one $i\in\left\{  1,2,\ldots,k\right\}  $ satisfies $\beta_{i}\neq0$ (since
the coefficients $\beta_{1},\beta_{2},\ldots,\beta_{k}$ are not all zero).
Pick the \textbf{largest} such $i$. Thus, $\beta_{i}\neq0$ but $\beta
_{i+1}=\beta_{i+2}=\cdots=\beta_{k}=0$. Hence,
\begin{align*}
&  \beta_{1}v_{1}+\beta_{2}v_{2}+\cdots+\beta_{k}v_{k}\\
&  =\left(  \beta_{1}v_{1}+\beta_{2}v_{2}+\cdots+\beta_{i-1}v_{i-1}\right)
+\beta_{i}v_{i}+\underbrace{\left(  \beta_{i+1}v_{i+1}+\beta_{i+2}%
v_{i+2}+\cdots+\beta_{k}v_{k}\right)  }_{\substack{=0v_{i+1}+0v_{i+2}%
+\cdots+0v_{k}\\=\mathbf{0}}}\\
&  =\left(  \beta_{1}v_{1}+\beta_{2}v_{2}+\cdots+\beta_{i-1}v_{i-1}\right)
+\beta_{i}v_{i},
\end{align*}
so that%
\begin{align*}
\beta_{i}v_{i}  &  =\underbrace{\left(  \beta_{1}v_{1}+\beta_{2}v_{2}%
+\cdots+\beta_{k}v_{k}\right)  }_{=\mathbf{0}}-\left(  \beta_{1}v_{1}%
+\beta_{2}v_{2}+\cdots+\beta_{i-1}v_{i-1}\right) \\
&  =-\left(  \beta_{1}v_{1}+\beta_{2}v_{2}+\cdots+\beta_{i-1}v_{i-1}\right) \\
&  =\left(  -\beta_{1}\right)  v_{1}+\left(  -\beta_{2}\right)  v_{2}%
+\cdots+\left(  -\beta_{i-1}\right)  v_{i-1}\\
&  \in\operatorname*{span}\left(  v_{1},v_{2},\ldots,v_{i-1}\right)  .
\end{align*}
Since $\beta_{i}\neq0$, we thus obtain $v_{i}\in\operatorname*{span}\left(
v_{1},v_{2},\ldots,v_{i-1}\right)  $ (since the set $\operatorname*{span}%
\left(  v_{1},v_{2},\ldots,v_{i-1}\right)  $ is an $F$-vector subspace of $V$
and thus preserved under scaling). This contradicts our assumption that
$v_{i}\notin\operatorname*{span}\left(  v_{1},v_{2},\ldots,v_{i-1}\right)  $.
This contradiction shows that our assumption was wrong, and thus completes our
proof of the \textquotedblleft$\Longleftarrow$\textquotedblright\ direction of
our claim.
\end{enumerate}

Thus, both directions of our claim are proved. This concludes the proof of
Lemma \ref{lem.linalg.lin-ind-via-span}.
\end{proof}

The next lemma we are going to use is itself an answer to a rather natural
counting problem. Indeed, it is well-known that (see, e.g., \cite[Proposition
2.7.2]{19fco}) if $X$ is an $n$-element set, and if $k\in\mathbb{N}$, then
the
\begin{align}
&  \left(  \text{\# of }k\text{-tuples of distinct elements of }X\right)
\nonumber\\
&  =n\left(  n-1\right)  \left(  n-2\right)  \cdots\left(  n-k+1\right)
=\prod_{i=0}^{k-1}\left(  n-i\right)  . \label{eq.count.k-tups-dist-elts}%
\end{align}
The following lemma is a \textquotedblleft linear analogue\textquotedblright%
\ of this combinatorial fact: The $n$-element set $X$ is replaced by an
$n$-dimensional vector space $V$, and \textquotedblleft distinct
elements\textquotedblright\ are replaced by \textquotedblleft linearly
independent vectors\textquotedblright. The answer is rather similar:

\begin{lemma}
\label{lem.pars.qbinom.lin-ind-count}Let $F$ be a finite field. Let
$n,k\in\mathbb{N}$. Let $V$ be an $n$-dimensional $F$-vector space. Then,%
\begin{align*}
&  \left(  \text{\# of linearly independent }k\text{-tuples of vectors in
}V\right) \\
&  =\left(  \left\vert F\right\vert ^{n}-\left\vert F\right\vert ^{0}\right)
\left(  \left\vert F\right\vert ^{n}-\left\vert F\right\vert ^{1}\right)
\cdots\left(  \left\vert F\right\vert ^{n}-\left\vert F\right\vert
^{k-1}\right)  =\prod_{i=0}^{k-1}\left(  \left\vert F\right\vert
^{n}-\left\vert F\right\vert ^{i}\right)  .
\end{align*}

\end{lemma}

\begin{proof}
[Proof of Lemma \ref{lem.pars.qbinom.lin-ind-count}.]We have $\left\vert
V\right\vert =\left\vert F\right\vert ^{n}$ (since $V$ is an $n$-dimensional
$F$-vector space).

Lemma \ref{lem.linalg.lin-ind-via-span} says that a $k$-tuple $\left(
v_{1},v_{2},\ldots,v_{k}\right)  $ of vectors in $V$ is linearly independent
if and only if it satisfies%
\begin{align*}
v_{1}  &  \notin\underbrace{\operatorname*{span}\left(  {}\right)
}_{=\left\{  \mathbf{0}\right\}  }\ \ \ \ \ \ \ \ \ \ \text{and}\\
v_{2}  &  \notin\operatorname*{span}\left(  v_{1}\right)
\ \ \ \ \ \ \ \ \ \ \text{and}\\
v_{3}  &  \notin\operatorname*{span}\left(  v_{1},v_{2}\right)
\ \ \ \ \ \ \ \ \ \ \text{and}\\
&  \cdots\ \ \ \ \ \ \ \ \ \ \text{and}\\
v_{k}  &  \notin\operatorname*{span}\left(  v_{1},v_{2},\ldots,v_{k-1}\right)
.
\end{align*}


Thus, we can construct a linearly independent $k$-tuple $\left(  v_{1}%
,v_{2},\ldots,v_{k}\right)  $ of vectors in $V$ as follows, proceeding entry
by entry:

\begin{itemize}
\item First, we choose $v_{1}$. This has to be a vector in $V\setminus
\underbrace{\operatorname*{span}\left(  {}\right)  }_{=\left\{  \mathbf{0}%
\right\}  }$ (because it has to satisfy $v_{1}\notin\operatorname*{span}%
\left(  {}\right)  $); thus, there are $\left\vert V\setminus
\operatorname*{span}\left(  {}\right)  \right\vert =\left\vert V\right\vert
-\underbrace{\left\vert \operatorname*{span}\left(  {}\right)  \right\vert
}_{=1}=\left\vert V\right\vert -1$ options for it.

Once $v_{1}$ has been chosen, we have obtained a linearly independent
singleton list $\left(  v_{1}\right)  $. Hence, its span $\operatorname*{span}%
\left(  v_{1}\right)  $ has dimension $1$ (as an $F$-vector space) and thus
size $\left\vert F\right\vert ^{1}$. In other words, $\left\vert
\operatorname*{span}\left(  v_{1}\right)  \right\vert =\left\vert F\right\vert
^{1}$.

\item Next, we choose $v_{2}$. This has to be a vector in $V\setminus
\operatorname*{span}\left(  v_{1}\right)  $ (because it has to satisfy
$v_{2}\notin\operatorname*{span}\left(  v_{1}\right)  $); thus, there are
$\left\vert V\right\vert -\underbrace{\left\vert \operatorname*{span}\left(
v_{1}\right)  \right\vert }_{=\left\vert F\right\vert ^{1}}=\left\vert
V\right\vert -\left\vert F\right\vert ^{1}$ options for it.

Once $v_{2}$ has been chosen, we have obtained a linearly independent list
$\left(  v_{1},v_{2}\right)  $. Hence, its span $\operatorname*{span}\left(
v_{1},v_{2}\right)  $ has dimension $2$ (as an $F$-vector space) and thus size
$\left\vert F\right\vert ^{2}$. In other words, $\left\vert
\operatorname*{span}\left(  v_{1},v_{2}\right)  \right\vert =\left\vert
F\right\vert ^{2}$.

\item Next, we choose $v_{3}$. This has to be a vector in $V\setminus
\operatorname*{span}\left(  v_{1},v_{2}\right)  $ (because it has to satisfy
$v_{3}\notin\operatorname*{span}\left(  v_{1},v_{2}\right)  $); thus, there
are $\left\vert V\right\vert -\underbrace{\left\vert \operatorname*{span}%
\left(  v_{1},v_{2}\right)  \right\vert }_{=\left\vert F\right\vert ^{2}%
}=\left\vert V\right\vert -\left\vert F\right\vert ^{2}$ options for it.

Once $v_{3}$ has been chosen, we have obtained a linearly independent list
$\left(  v_{1},v_{2},v_{3}\right)  $. Hence, its span $\operatorname*{span}%
\left(  v_{1},v_{2},v_{3}\right)  $ has dimension $3$ (as an $F$-vector space)
and thus size $\left\vert F\right\vert ^{3}$. In other words, $\left\vert
\operatorname*{span}\left(  v_{1},v_{2},v_{3}\right)  \right\vert =\left\vert
F\right\vert ^{3}$.

\item And so on, until the last vector $v_{k}$ in our list has been chosen.
\end{itemize}

The total \# of ways to perform this construction is%
\[
\left(  \left\vert V\right\vert -1\right)  \left(  \left\vert V\right\vert
-\left\vert F\right\vert ^{1}\right)  \left(  \left\vert V\right\vert
-\left\vert F\right\vert ^{2}\right)  \cdots\left(  \left\vert V\right\vert
-\left\vert F\right\vert ^{k-1}\right)  .
\]


Hence,
\begin{align*}
&  \left(  \text{\# of linearly independent }k\text{-tuples of vectors in
}V\right) \\
&  =\left(  \left\vert V\right\vert -1\right)  \left(  \left\vert V\right\vert
-\left\vert F\right\vert ^{1}\right)  \left(  \left\vert V\right\vert
-\left\vert F\right\vert ^{2}\right)  \cdots\left(  \left\vert V\right\vert
-\left\vert F\right\vert ^{k-1}\right) \\
&  =\prod_{i=0}^{k-1}\left(  \left\vert V\right\vert -\left\vert F\right\vert
^{i}\right)  =\prod_{i=0}^{k-1}\left(  \left\vert F\right\vert ^{n}-\left\vert
F\right\vert ^{i}\right)
\end{align*}
(since $\left\vert V\right\vert =\left\vert F\right\vert ^{n}$). This proves
Lemma \ref{lem.pars.qbinom.lin-ind-count}.
\end{proof}

Another lemma we will need is a basic combinatorial principle (often
illustrated by the saying \textquotedblleft to count a flock of sheep, count
the legs and divide by $4$\textquotedblright):

\begin{lemma}
[Multijection principle]\label{lem.count.multijection}Let $A$ and $B$ be two
finite sets. Let $m\in\mathbb{N}$. Let $f:A\rightarrow B$ be any map. Assume
that each $b\in B$ has exactly $m$ preimages under $f$ (that is, for each
$b\in B$, there are exactly $m$ many elements $a\in A$ such that $f\left(
a\right)  =b$). Then,%
\[
\left\vert A\right\vert =m\cdot\left\vert B\right\vert .
\]

\end{lemma}

\begin{proof}
Easy and LTTR.
\end{proof}

We note that a map $f:A\rightarrow B$ satisfying the assumption of Lemma
\ref{lem.count.multijection} is often called an $m$\emph{-to-}$1$\emph{ map}.

\begin{proof}
[Proof of Theorem \ref{thm.pars.qbinom.subsp-count}.]First of all, we notice
that if $k>n$, then $\dbinom{n}{k}_{\left\vert F\right\vert }=0$ (by
Proposition \ref{prop.pars.qbinom.0}) and $\left(  \text{\# of }%
k\text{-dimensional vector subspaces of }V\right)  =0$ (since the dimension of
a subspace of $V$ is never larger than the dimension of $V$). Thus, Theorem
\ref{thm.pars.qbinom.subsp-count} is true when $k>n$. Hence, for the rest of
this proof, we WLOG assume that $k\leq n$.

We will use the shorthand \textquotedblleft linind\textquotedblright\ for the
words \textquotedblleft linearly independent\textquotedblright.

If $\left(  v_{1},v_{2},\ldots,v_{k}\right)  $ is a linind $k$-tuple of
vectors in $V$, then $\operatorname*{span}\left(  v_{1},v_{2},\ldots
,v_{k}\right)  $ is a $k$-dimensional vector subspace of $V$. Hence, we can
define a map%
\begin{align*}
f:\left\{  \text{linind }k\text{-tuples of vectors in }V\right\}   &
\rightarrow\left\{  k\text{-dimensional vector subspaces of }V\right\}  ,\\
\left(  v_{1},v_{2},\ldots,v_{k}\right)   &  \mapsto\operatorname*{span}%
\left(  v_{1},v_{2},\ldots,v_{k}\right)  .
\end{align*}
Consider this map $f$. We claim the following:

\begin{statement}
\textit{Observation 1:} Each $k$-dimensional vector subspace of $V$ has
exactly%
\[
\left(  \left\vert F\right\vert ^{k}-\left\vert F\right\vert ^{0}\right)
\left(  \left\vert F\right\vert ^{k}-\left\vert F\right\vert ^{1}\right)
\cdots\left(  \left\vert F\right\vert ^{k}-\left\vert F\right\vert
^{k-1}\right)
\]
preimages under $f$.
\end{statement}

[\textit{Proof of Observation 1:} Let $W$ be a $k$-dimensional vector subspace
of $V$. We must prove that
\[
\left(  \text{\# of preimages of }W\text{ under }f\right)  =\left(  \left\vert
F\right\vert ^{k}-\left\vert F\right\vert ^{0}\right)  \left(  \left\vert
F\right\vert ^{k}-\left\vert F\right\vert ^{1}\right)  \cdots\left(
\left\vert F\right\vert ^{k}-\left\vert F\right\vert ^{k-1}\right)  .
\]


A preimage of $W$ under $f$ is a $k$-tuple of vectors in $V$ that spans $W$
(by the very definition of $f$). Obviously, all the $k$ vectors in such a
$k$-tuple must belong to $W$. Thus, a preimage of $W$ under $f$ is a $k$-tuple
of vectors in $W$ that spans $W$. However, the vector space $W$ is
$k$-dimensional. Thus, a $k$-tuple of vectors in $W$ spans $W$ if and only if
this $k$-tuple is linind\footnote{Again, we are using a simple fact from
linear algebra here, which is true over any field (not necessarily finite): A
$k$-tuple of vectors in a $k$-dimensional vector space spans the space if and
only if it is linind.}. Therefore, a preimage of $W$ under $f$ is a linind
$k$-tuple of vectors in $W$. Hence,%
\begin{align*}
&  \left(  \text{\# of preimages of }W\text{ under }f\right) \\
&  =\left(  \text{\# of linearly independent }k\text{-tuples of vectors in
}W\right) \\
&  =\left(  \left\vert F\right\vert ^{k}-\left\vert F\right\vert ^{0}\right)
\left(  \left\vert F\right\vert ^{k}-\left\vert F\right\vert ^{1}\right)
\cdots\left(  \left\vert F\right\vert ^{k}-\left\vert F\right\vert
^{k-1}\right)
\end{align*}
(by Lemma \ref{lem.pars.qbinom.lin-ind-count}, applied to $W$ and $k$ instead
of $V$ and $n$). This proves Observation 1.]

Now, Observation 1 shows that each $k$-dimensional vector subspace of $V$ has
exactly $\left(  \left\vert F\right\vert ^{k}-\left\vert F\right\vert
^{0}\right)  \left(  \left\vert F\right\vert ^{k}-\left\vert F\right\vert
^{1}\right)  \cdots\left(  \left\vert F\right\vert ^{k}-\left\vert
F\right\vert ^{k-1}\right)  $ preimages under $f$. Hence, Lemma
\ref{lem.count.multijection} (applied to $A=\left\{  \text{linind
}k\text{-tuples of vectors in }V\right\}  $ and $B=\left\{
k\text{-dimensional vector subspaces of }V\right\}  $ and \newline$m=\left(
\left\vert F\right\vert ^{k}-\left\vert F\right\vert ^{0}\right)  \left(
\left\vert F\right\vert ^{k}-\left\vert F\right\vert ^{1}\right)
\cdots\left(  \left\vert F\right\vert ^{k}-\left\vert F\right\vert
^{k-1}\right)  $) shows that%
\begin{align*}
&  \left(  \text{\# of linind }k\text{-tuples of vectors in }V\right) \\
&  =\left(  \left\vert F\right\vert ^{k}-\left\vert F\right\vert ^{0}\right)
\left(  \left\vert F\right\vert ^{k}-\left\vert F\right\vert ^{1}\right)
\cdots\left(  \left\vert F\right\vert ^{k}-\left\vert F\right\vert
^{k-1}\right) \\
&  \ \ \ \ \ \ \ \ \ \ \cdot\left(  \text{\# of }k\text{-dimensional vector
subspaces of }V\right)  .
\end{align*}
However, Lemma \ref{lem.pars.qbinom.lin-ind-count} yields%
\begin{align*}
&  \left(  \text{\# of linind }k\text{-tuples of vectors in }V\right) \\
&  =\left(  \left\vert F\right\vert ^{n}-\left\vert F\right\vert ^{0}\right)
\left(  \left\vert F\right\vert ^{n}-\left\vert F\right\vert ^{1}\right)
\cdots\left(  \left\vert F\right\vert ^{n}-\left\vert F\right\vert
^{k-1}\right)  .
\end{align*}
Comparing these two equalities, we obtain%
\begin{align*}
&  \left(  \left\vert F\right\vert ^{k}-\left\vert F\right\vert ^{0}\right)
\left(  \left\vert F\right\vert ^{k}-\left\vert F\right\vert ^{1}\right)
\cdots\left(  \left\vert F\right\vert ^{k}-\left\vert F\right\vert
^{k-1}\right) \\
&  \ \ \ \ \ \ \ \ \ \ \cdot\left(  \text{\# of }k\text{-dimensional vector
subspaces of }V\right) \\
&  =\left(  \left\vert F\right\vert ^{n}-\left\vert F\right\vert ^{0}\right)
\left(  \left\vert F\right\vert ^{n}-\left\vert F\right\vert ^{1}\right)
\cdots\left(  \left\vert F\right\vert ^{n}-\left\vert F\right\vert
^{k-1}\right)  .
\end{align*}
Therefore,%
\begin{align*}
&  \left(  \text{\# of }k\text{-dimensional vector subspaces of }V\right) \\
&  =\dfrac{\left(  \left\vert F\right\vert ^{n}-\left\vert F\right\vert
^{0}\right)  \left(  \left\vert F\right\vert ^{n}-\left\vert F\right\vert
^{1}\right)  \cdots\left(  \left\vert F\right\vert ^{n}-\left\vert
F\right\vert ^{k-1}\right)  }{\left(  \left\vert F\right\vert ^{k}-\left\vert
F\right\vert ^{0}\right)  \left(  \left\vert F\right\vert ^{k}-\left\vert
F\right\vert ^{1}\right)  \cdots\left(  \left\vert F\right\vert ^{k}%
-\left\vert F\right\vert ^{k-1}\right)  }\\
&  =\dfrac{\prod_{i=0}^{k-1}\left(  \left\vert F\right\vert ^{n}-\left\vert
F\right\vert ^{i}\right)  }{\prod_{i=0}^{k-1}\left(  \left\vert F\right\vert
^{k}-\left\vert F\right\vert ^{i}\right)  }=\prod_{i=0}^{k-1}%
\underbrace{\dfrac{\left\vert F\right\vert ^{n}-\left\vert F\right\vert ^{i}%
}{\left\vert F\right\vert ^{k}-\left\vert F\right\vert ^{i}}}%
_{\substack{=\dfrac{\left\vert F\right\vert ^{i}\left(  \left\vert
F\right\vert ^{n-i}-1\right)  }{\left\vert F\right\vert ^{i}\left(  \left\vert
F\right\vert ^{k-i}-1\right)  }=\dfrac{\left\vert F\right\vert ^{n-i}%
-1}{\left\vert F\right\vert ^{k-i}-1}\\=\dfrac{1-\left\vert F\right\vert
^{n-i}}{1-\left\vert F\right\vert ^{k-i}}}}=\prod_{i=0}^{k-1}\dfrac
{1-\left\vert F\right\vert ^{n-i}}{1-\left\vert F\right\vert ^{k-i}}\\
&  =\dfrac{\prod_{i=0}^{k-1}\left(  1-\left\vert F\right\vert ^{n-i}\right)
}{\prod_{i=0}^{k-1}\left(  1-\left\vert F\right\vert ^{k-i}\right)  }%
=\dfrac{\left(  1-\left\vert F\right\vert ^{n}\right)  \left(  1-\left\vert
F\right\vert ^{n-1}\right)  \cdots\left(  1-\left\vert F\right\vert
^{n-k+1}\right)  }{\left(  1-\left\vert F\right\vert ^{k}\right)  \left(
1-\left\vert F\right\vert ^{k-1}\right)  \cdots\left(  1-\left\vert
F\right\vert ^{1}\right)  }\\
&  =\dbinom{n}{k}_{\left\vert F\right\vert }%
\end{align*}
(since substituting $\left\vert F\right\vert $ for $q$ in Theorem
\ref{thm.pars.qbinom.quot1} \textbf{(b)} yields \newline$\dbinom{n}%
{k}_{\left\vert F\right\vert }=\dfrac{\left(  1-\left\vert F\right\vert
^{n}\right)  \left(  1-\left\vert F\right\vert ^{n-1}\right)  \cdots\left(
1-\left\vert F\right\vert ^{n-k+1}\right)  }{\left(  1-\left\vert F\right\vert
^{k}\right)  \left(  1-\left\vert F\right\vert ^{k-1}\right)  \cdots\left(
1-\left\vert F\right\vert ^{1}\right)  }$). This proves Theorem
\ref{thm.pars.qbinom.subsp-count}.
\end{proof}

\subsubsection{Limits of $q$-binomial coefficients}

There is much more to say about $q$-binomial coefficients, but let us just
briefly focus on their limiting behavior. This is not analogous to anything
known from usual binomial coefficients; indeed, the limit $\lim
\limits_{n\rightarrow\infty}\dbinom{n}{k}$ does not exist for any positive
integer $k$. However, $q$-binomial coefficients behave much better in this regard.

Indeed, consider the $q$-binomial coefficients $\dbinom{n}{2}_{q}$ for various
values of $n$:%
\begin{align*}
\dbinom{0}{2}_{q}  &  =0,\\
\dbinom{1}{2}_{q}  &  =0,\\
\dbinom{2}{2}_{q}  &  =1,\\
\dbinom{3}{2}_{q}  &  =1+q+q^{2},\\
\dbinom{4}{2}_{q}  &  =1+q+2q^{2}+q^{3}+q^{4},\\
\dbinom{5}{2}_{q}  &  =1+q+2q^{2}+2q^{3}+2q^{4}+q^{5}+q^{6},\\
\dbinom{6}{2}_{q}  &  =1+q+2q^{2}+2q^{3}+3q^{4}+2q^{5}+2q^{6}+q^{7}+q^{8}.
\end{align*}
It appears from these examples that the sequence $\left(  \dbinom{n}{2}%
_{q}\right)  _{n\in\mathbb{N}}$ coefficientwise stabilizes\footnote{Recall
Definition \ref{def.fps.lim.coeff-stab} for the notion of \textquotedblleft
coefficientwise stabilizing\textquotedblright.} to%
\[
1+q+2q^{2}+2q^{3}+3q^{4}+3q^{5}+\cdots=\sum_{n\in\mathbb{N}}\left(
1+\left\lfloor \dfrac{n}{2}\right\rfloor \right)  q^{n}.
\]
And this is indeed the case:

\begin{proposition}
\label{prop.pars.qbinom.lim1}Let $k\in\mathbb{N}$ be fixed. Then,%
\[
\lim\limits_{n\rightarrow\infty}\dbinom{n}{k}_{q}=\sum_{n\in\mathbb{N}}\left(
p_{0}\left(  n\right)  +p_{1}\left(  n\right)  +\cdots+p_{k}\left(  n\right)
\right)  q^{n}=\prod_{i=1}^{k}\dfrac{1}{1-q^{i}}.
\]
(See Definition \ref{def.pars.pn-pkn} \textbf{(a)} for the meaning of
$p_{i}\left(  n\right)  $.)
\end{proposition}

\begin{proof}
[First proof of Proposition \ref{prop.pars.qbinom.lim1} (sketched).]For each
integer $n\geq k$, we have%
\begin{align*}
\dbinom{n}{k}_{q}  &  =\dfrac{\left(  1-q^{n}\right)  \left(  1-q^{n-1}%
\right)  \cdots\left(  1-q^{n-k+1}\right)  }{\left(  1-q^{k}\right)  \left(
1-q^{k-1}\right)  \cdots\left(  1-q^{1}\right)  }\ \ \ \ \ \ \ \ \ \ \left(
\text{by Theorem \ref{thm.pars.qbinom.quot1} \textbf{(b)}}\right) \\
&  =\dfrac{\left(  1-q^{n-k+1}\right)  \left(  1-q^{n-k+2}\right)
\cdots\left(  1-q^{n}\right)  }{\left(  1-q^{1}\right)  \left(  1-q^{2}%
\right)  \cdots\left(  1-q^{k}\right)  }\\
&  \ \ \ \ \ \ \ \ \ \ \ \ \ \ \ \ \ \ \ \ \left(  \text{here, we have turned
both products upside down}\right) \\
&  =\dfrac{\prod_{i=1}^{k}\left(  1-q^{n-k+i}\right)  }{\prod_{i=1}^{k}\left(
1-q^{i}\right)  }=\dfrac{1}{\prod_{i=1}^{k}\left(  1-q^{i}\right)  }\cdot
\prod_{i=1}^{k}\left(  1-q^{n-k+i}\right)  .
\end{align*}


However, we have $\lim\limits_{n\rightarrow\infty}q^{n}=0$ (check
this!\footnote{This is a matter of understanding Definition
\ref{def.fps.lim.coeff-stab}.}). Thus, for each $i\in\left\{  1,2,\ldots
,k\right\}  $, we have%
\[
\lim\limits_{n\rightarrow\infty}q^{n-k+i}=0
\]
(since the family $\left(  q^{n-k+i}\right)  _{n\geq k-i}$ is just a
reindexing of the family $\left(  q^{n}\right)  _{n\geq0}$), and therefore%
\[
\lim\limits_{n\rightarrow\infty}\left(  1-q^{n-k+i}\right)  =1.
\]
Hence, Corollary \ref{cor.fps.lim.sum-prod-k} (applied to $f_{i,n}%
=1-q^{n-k+i}$ and $f_{i}=1$) yields that%
\[
\lim\limits_{n\rightarrow\infty}\sum_{i=1}^{k}\left(  1-q^{n-k+i}\right)
=\sum_{i=1}^{k}1\ \ \ \ \ \ \ \ \ \ \text{and}\ \ \ \ \ \ \ \ \ \ \lim
\limits_{n\rightarrow\infty}\prod_{i=1}^{k}\left(  1-q^{n-k+i}\right)
=\prod_{i=1}^{k}1.
\]
Thus, in particular,%
\[
\lim\limits_{n\rightarrow\infty}\prod_{i=1}^{k}\left(  1-q^{n-k+i}\right)
=\prod_{i=1}^{k}1=1.
\]


Now, recall that each integer $n\geq k$ satisfies%
\[
\dbinom{n}{k}_{q}=\dfrac{1}{\prod_{i=1}^{k}\left(  1-q^{i}\right)  }\cdot
\prod_{i=1}^{k}\left(  1-q^{n-k+i}\right)  .
\]
Hence,%
\begin{align}
\lim\limits_{n\rightarrow\infty}\dbinom{n}{k}_{q}  &  =\lim
\limits_{n\rightarrow\infty}\left(  \dfrac{1}{\prod_{i=1}^{k}\left(
1-q^{i}\right)  }\cdot\prod_{i=1}^{k}\left(  1-q^{n-k+i}\right)  \right)
\nonumber\\
&  =\dfrac{1}{\prod_{i=1}^{k}\left(  1-q^{i}\right)  }\cdot\underbrace{\lim
\limits_{n\rightarrow\infty}\prod_{i=1}^{k}\left(  1-q^{n-k+i}\right)  }%
_{=1}=\dfrac{1}{\prod_{i=1}^{k}\left(  1-q^{i}\right)  }\nonumber\\
&  =\prod_{i=1}^{k}\dfrac{1}{1-q^{i}}. \label{pf.prop.pars.qbinom.lim1.rhs1}%
\end{align}


Finally, Theorem \ref{thm.pars.main-gf-0n} (with the letters $x$, $m$ and $k$
renamed as $q$, $k$ and $i$) says that%
\[
\sum_{n\in\mathbb{N}}\left(  p_{0}\left(  n\right)  +p_{1}\left(  n\right)
+\cdots+p_{k}\left(  n\right)  \right)  q^{n}=\prod_{i=1}^{k}\dfrac{1}%
{1-q^{i}}.
\]
Combining this with (\ref{pf.prop.pars.qbinom.lim1.rhs1}), we obtain%
\[
\lim\limits_{n\rightarrow\infty}\dbinom{n}{k}_{q}=\sum_{n\in\mathbb{N}}\left(
p_{0}\left(  n\right)  +p_{1}\left(  n\right)  +\cdots+p_{k}\left(  n\right)
\right)  q^{n}=\prod_{i=1}^{k}\dfrac{1}{1-q^{i}}.
\]
This proves Proposition \ref{prop.pars.qbinom.lim1}.
\end{proof}

\begin{proof}
[Second proof of Proposition \ref{prop.pars.qbinom.lim1} (sketched).]For each
$n\in\mathbb{N}$, we have%
\begin{equation}
\dbinom{n}{k}_{q}=\sum_{\substack{\lambda\text{ is a partition}\\\text{with
largest part }\leq n-k\\\text{and length }\leq k}}q^{\left\vert \lambda
\right\vert } \label{pf.prop.pars.qbinom.lim1.2nd.1}%
\end{equation}
(by the definition of $\dbinom{n}{k}_{q}$).

However, for each $n\in\mathbb{N}$, the sum $\sum_{\substack{\lambda\text{ is
a partition}\\\text{with largest part }\leq n-k\\\text{and length }\leq
k}}q^{\left\vert \lambda\right\vert }$ is a partial sum of the sum
$\sum_{\substack{\lambda\text{ is a partition}\\\text{with length }\leq
k}}q^{\left\vert \lambda\right\vert }$, and this partial sum grows by more and
more addends as $n$ increases; each addend of the sum $\sum_{\substack{\lambda
\text{ is a partition}\\\text{with length }\leq k}}q^{\left\vert
\lambda\right\vert }$ gets eventually included in this partial sum (for
sufficiently large $n$). From these observations, it is easy to obtain that%
\[
\lim\limits_{n\rightarrow\infty}\sum_{\substack{\lambda\text{ is a
partition}\\\text{with largest part }\leq n-k\\\text{and length }\leq
k}}q^{\left\vert \lambda\right\vert }=\sum_{\substack{\lambda\text{ is a
partition}\\\text{with length }\leq k}}q^{\left\vert \lambda\right\vert }.
\]
In view of (\ref{pf.prop.pars.qbinom.lim1.2nd.1}), this rewrites as%
\begin{align*}
\lim\limits_{n\rightarrow\infty}\dbinom{n}{k}_{q}  &  =\sum_{\substack{\lambda
\text{ is a partition}\\\text{with length }\leq k}}q^{\left\vert
\lambda\right\vert }=\sum_{n\in\mathbb{N}}\underbrace{\left(  \text{\# of
partitions of }n\text{ having length }\leq k\right)  }_{\substack{=p_{0}%
\left(  n\right)  +p_{1}\left(  n\right)  +\cdots+p_{k}\left(  n\right)
\\\text{(by Definition \ref{def.pars.pn-pkn} \textbf{(a)})}}}q^{n}\\
&  =\sum_{n\in\mathbb{N}}\left(  p_{0}\left(  n\right)  +p_{1}\left(
n\right)  +\cdots+p_{k}\left(  n\right)  \right)  q^{n}=\prod_{i=1}^{k}%
\dfrac{1}{1-q^{i}}%
\end{align*}
(by Theorem \ref{thm.pars.main-gf-0n}, with the letters $x$, $m$ and $k$
renamed as $q$, $k$ and $i$). Thus, Proposition \ref{prop.pars.qbinom.lim1} is
proved again.
\end{proof}
